% !TeX program = xelatex
\documentclass[11pt,a4paper]{article}
\usepackage[margin=23mm]{geometry}
\usepackage{fontspec}
\setmainfont{Latin Modern Roman}
\usepackage{amsmath,amssymb,mathtools}
\usepackage{xcolor,enumitem,booktabs,longtable,array,xurl,hyperref}
\definecolor{muted}{HTML}{53636C}
\hypersetup{colorlinks=true,urlcolor=blue,linkcolor=blue}
\setlength{\parindent}{0pt}
\setlength{\parskip}{5pt}
\newcommand{\R}{\mathbb R}
\newcommand{\E}{\mathbb E}
\newcommand{\N}{\mathbb N}
\newcommand{\ind}{\mathbf 1}
\newcommand{\Var}{\operatorname{Var}}
\newcommand{\Cov}{\operatorname{Cov}}
\newcommand{\supp}{\operatorname{supp}}
\DeclareMathOperator*{\argmin}{arg\,min}
\DeclareMathOperator*{\argmax}{arg\,max}
\newcommand{\entrymeta}[3]{{\sffamily\small\color{muted}\textbf{\texttt{#1}}\quad|\quad #2\par #3\par}}
\newcommand{\Problemref}[1]{\texttt{#1}}
\newcommand{\JELref}[2]{\texttt{#2} (JEL \texttt{#1})}
\begin{document}
\section*{Policy effects beyond IV compliers}
\textbf{Problem ID:} \texttt{C26-403}\quad \textbf{JEL:} \texttt{C26}\par
% BEGIN ECONOMICS STATEMENT
\label{problem:OP-0424}
{\sffamily\small\color{muted}Original subject group: Econometrics, causal inference and measurement\par}
\label{definitions:problem:30007706}
\textbf{Audit classification:} Background; proposed extension. The verified chapter surveys solutions and recommends target-specific modeling; it does not identify this model class as a single open theorem.
\subsection*{Definitions and precise research target}
Consider a population with covariates \(X\), binary treatment \(D\), bounded potential outcomes \(Y(0),Y(1)\in[0,1]\), and a possibly multivalued instrument \(Z\). Latent vector \(V\) describes treatment choices and effect heterogeneity. Assume consistency, instrument exclusion, and conditional independence of \(Z\) from potential outcomes and potential treatment choices given \(X\); impose monotonicity only where justified. For each model \(M\) satisfying these restrictions, let \(D_M^q(V,X)\) denote treatment under a specified new assignment policy \(q\), and \(D_M^b(V,X)\) treatment under the baseline policy \(b\). The policy effect is
\[
\psi(M)=E_M[Y(D_M^q(V,X))-Y(D_M^b(V,X))].
\]
Define \(\mathcal M(P,\mathcal R)\) as models reproducing the observed law \(P\) of \((Y,D,Z,X)\), where \(Y=Y(D)\), and obeying additional restrictions \(\mathcal R\), and define the identified set \(\{\psi(M):M\in\mathcal M(P,\mathcal R)\}\). Which credible structural restrictions, new instruments, randomized policy pilots or outcome bounds narrow this set enough to evaluate \(q\)? Clearly specify how each model maps the proposed policy into treatment choices, rather than equating a complier effect with the policy effect. Existing IV and partial-identification methods already solve important cases; this proposed problem asks which added information is credible in an application with multidimensional heterogeneity.

The application, target population, outcome measurement and horizon, concrete policies \(q,b\), and criterion for ``enough to evaluate'' remain to be specified; no numerical informativeness threshold or decision rule is selected here. Each concrete variant must declare \(\mathcal R\). Any monotonicity restriction must name the instrument-state comparisons and the direction of the almost-sure treatment inequalities conditional on \(X\); no ordering of a vector instrument is implicit. Added instruments or pilots require an augmented observable law and explicit links to the target population. A narrowing claim requires that models compatible with the added information project into the original compatible class while preserving \(\psi\); incompatibility must be reported separately.

\textbf{Definitions, assumptions and mathematical qualifications.} A model supplies \(Y_M(d,V,X)\), an instrument-to-choice law \(D_M(z,V,X)\), the joint law of \((V,X,Z)\), including the conditional instrument assignment law, and explicit policy maps for \(q,b\). It specifies measurable domains and measurable, binary-valued treatment maps. Baseline treatment is \(D_M(Z,V,X)\). In the displayed target, \(D_M^b(V,X)\) is shorthand for this baseline random treatment; \(D_M^q(V,X)\) similarly denotes the policy treatment on the model's probability space. The shorthand does not assert that realized baseline treatment is deterministic given latent traits and covariates alone. For randomized policies, the model must supply the assignment randomness and its joint law with the other primitives; their dependence is not selected here. Full instrument exogeneity is \(Z\perp\{Y_M(0),Y_M(1),D_M(z):z\in\mathcal Z\}\mid X\); exclusion removes \(Z\) from outcomes apart from \(D\). Since outcomes are bounded, \(\psi(M)\in[-1,1]\). A nonempty observationally compatible class yields the exact attainable identified set. Its interval structure requires a separate argument, for example a mixture operation preserving compatibility and restrictions under which \(\psi\) is affine. Closedness alone does not imply interval structure; endpoint attainment is a separate question. Policy consistency and stable treatment versions must hold before extrapolating from compliers. The displayed target compares treatment assignments in a fixed population using the same treatment-specific outcomes and latent-trait distribution. Instrument exclusion alone does not justify this policy invariance. Whether the proposed intervention satisfies it remains application-specific; direct policy effects, changing treatment versions or spillovers are not automatically represented by this formula.
\subsection*{Variants and assumption changes}
\begin{enumerate}
\item \textbf{Sourced threshold benchmark.} Use \(D=\mathbf1\{U\leq p(Z,X)\}\), \(U\mid X\sim U[0,1]\), with conditional instrument independence. Define \(p(z,x)=\Pr(D=1\mid Z=z,X=x)\) and \(\operatorname{MTE}(u,x)=E[Y(1)-Y(0)\mid U=u,X=x]\), defined almost everywhere. Policy integrals require specified policy selection and invariance. Local-IV identification uses continuous conditional propensity variation and suitable differentiability; finite support alone does not identify an unrestricted pointwise MTE, although supported interval integrals may be identified. Each policy-integral claim must justify its weights, support and any functional restrictions, including extrapolation.
\item \textbf{Editorial multidimensional policy.} Permit vector \(V\) and policies affecting costs or eligibility outside observed support. Use pilots, added instruments and explicit bounds on selection or outcomes to identify a decision-relevant range for \(\psi\).
\end{enumerate}
\subsection*{References and source audit}
\textbf{Checked source.} 10 December 2024 author chapter, Sections 2.6 (p. 14), 4.3 (pp. 57--76), 4.7 (pp. 88--93), and 6 (p. 99). The chapter surveys policy-relevant IV, MTE and bounding solutions. Full primary text retrieved. \url{https://www.rfberlin.com/wp-content/uploads/2024/12/24030.pdf}.
\begin{enumerate}\raggedright
\item\hypertarget{definitions:source:30007706:1}{} Magne Mogstad; Alexander Torgovitsky. 2024. \emph{Instrumental Variables with Unobserved Heterogeneity in Treatment Effects}. 10 December 2024 author chapter, Sections 2.6 (p. 14), 4.3 (pp. 57--76), 4.7 (pp. 88--93), and 6 (p. 99).
\url{https://www.rfberlin.com/wp-content/uploads/2024/12/24030.pdf}.
DOI: \url{https://doi.org/10.3386/w32927}.
\end{enumerate}

\subsection*{Common conventions: Source mathematical conventions}

These conventions apply to each entry unless explicitly replaced.

\textbf{Models and identification.} A primitive model \(m\) specifies all probability laws, feasible choices, information, equilibrium and measurement rules used by its target. The admissible class \(\mathcal M\) is fixed independently of outcomes. If \(P_O(m)\) is its observable probability law and \(\tau(m)\) the target, its identified set at observable law \(P\) is
\[
 \mathcal I_\tau(P)=\{\tau(m):m\in\mathcal M,\ P_O(m)=P\}.
\]
Point identification means that this set is a singleton. An empty set signals incompatibility of data and assumptions. Coordinate bounds need not describe the joint set. The exact identified set contains attainable targets; its closure also contains limits of attainable targets. A reported set is exact only if it equals the attainable set, and is a sharp closed region only if it equals its closure. Sharp lower and upper endpoint bounds are the infimum and supremum of the attainable set and need not be attained. Finite bounds require explicit restrictions on unobserved quantities; unrestricted confounding, hidden wealth, extrapolation or arbitrary equilibrium selection can yield uninformative sets.

\textbf{Causal interventions.} A potential outcome \(Y_i(p)\) is the outcome under a complete policy regime \(p\), including timing, financing, information and market spillovers. The target population and units are fixed before treatment. With interference, use the complete assignment vector or a justified exposure mapping; individual no-interference notation alone cannot identify an economy-wide intervention. A difference of mean potential outcomes does not require the cross-world joint distribution. Conditioning on post-treatment migration, survival, enrollment or employment changes the estimand and needs separate justification.

\textbf{Statistical statements.} An estimator, confidence set, forecast or test specifies a sampling law, model class and loss/coverage criterion. Uniform \((1-\alpha)\)-coverage means \(\inf_{m\in\mathcal M}\Pr_m\{\tau(m)\in C_n\}\geq1-\alpha\), or an explicitly asymptotic version. The whole parameter space trivially has coverage, so informativeness requires a stated width, risk or power objective. A forecasting improvement is relative to a named benchmark, information set and evaluation population.

\textbf{Equilibrium and optimization.} An equilibrium comprises feasible plans, optimality, beliefs where needed, budget constraints and all market/resource clearing. The primitives or a stated benchmark define each correspondence \(\mathcal E(p;m)\). Nonemptiness is assumed or proved, never inferred just from compact policy sets. A selected-equilibrium target uses a declared selection rule; a robust target quantifies over all permitted equilibria. Use suprema or infima unless attainment is established. An \(\varepsilon\)-optimal policy is feasible and within \(\varepsilon\) of the finite optimum value. Report an empty feasible set separately.

\textbf{Welfare and accounting.} Normative utility, interpersonal normalization, social weights, numeraire, discounting and the use of public receipts are primitives. Behavioral utility need not coincide with the welfare criterion. Household welfare includes ownership income and rents once; adding profits again to compensating variation generally double-counts them. A monetary equivalent is defined by an explicit transfer and price convention, with monotonicity and range conditions. Average effects, mechanism contributions, transfers and welfare losses are different targets. Infinite sums require convergence and dynamic feasible plans require terminal or no-Ponzi conditions.

\textbf{Source versions.} A working-paper date, revision date and journal publication year are distinguished. A DOI for a journal version is not automatically the DOI of an earlier manuscript. Locators refer to the stated version and printed pages where specified. A metadata-only check does not verify a claim about a full-text conclusion. Every entry's source subsection narrows the provenance claim accordingly.
% END ECONOMICS STATEMENT
\end{document}
